\documentclass[11pt]{article}

\usepackage[margin=1in]{geometry}
\usepackage{amsmath,amssymb}
\usepackage{booktabs}
\usepackage{enumitem}
\usepackage{xurl}
\usepackage{hyperref}
\usepackage{xcolor}
\urlstyle{tt}

\newcommand{\method}{CAVE-Mem}
\newcommand{\score}{\mathrm{score}}
\newcommand{\valid}{\mathrm{valid}}
\newcommand{\boundary}{\mathrm{boundary}}
\newcommand{\utility}{\mathrm{utility}}

\title{
  When Should a Memory Agent Trust Its Experience?\\
  Conditional Applicability and Validated Experience for Memory Search
}
\author{Anonymous Authors}
\date{}

\begin{document}
\maketitle

\begin{abstract}
Experience-based memory agents reuse lessons distilled from previous search
trajectories, but most existing systems treat experience as unconditionally
helpful once it is semantically retrieved. This assumption is brittle:
reflections extracted from successful trajectories can be post-hoc,
domain-specific, or correct only under a particular memory substrate and answer
contract. We study this problem in long-context memory search and propose
\method, a training-free framework for conditional experience reuse.
\method{} represents an experience intervention by three additional fields
beyond its natural-language guidance: an applicability condition, a failure
boundary, and an observed utility estimate. At inference time, the agent first
profiles the memory substrate and the question's answer contract, then applies
candidate experience only when the condition holds and no boundary is crossed;
otherwise it abstains and falls back to the base memory-search agent. This turns
experience reuse from flat Top-$K$ prompt injection into a validated
intervention-selection problem. Using GPT-4o-mini as the backbone, \method{}
outperforms GAM and R2Mem on matched full/common splits of LoCoMo, HotpotQA,
and NarrativeQA, improving F1 by +5.85, +2.16, and +0.71 over GAM, and by
+6.47, +3.49, and +2.12 over R2Mem, respectively. The results suggest that the
central question for self-improving memory agents is not merely which
experience is relevant, but under which conditions that experience should be
trusted.
\end{abstract}

\section{Introduction}

Memory-search agents answer questions by repeatedly planning searches over a
stored memory, retrieving evidence, integrating the evidence into a working
summary, and producing an answer. Recent systems improve this loop by storing
experiences distilled from previous trajectories and injecting retrieved
experience into future searches. R2Mem is a representative example: it reflects
on good and bad search trajectories, builds an experience bank, and retrieves
Top-$K$ experience snippets during inference. This is appealing because it is
RL-free and can be added to existing agents.

However, flat experience reuse hides a stronger assumption: if an experience is
similar to the current state, it should be reused. In practice, the same
experience can help in one condition and hurt in another. A temporal search
rule can be useful for episodic dialogue but misleading for narrative plots. A
counting rule can help when raw mentions are recoverable, but hurt when the
question asks for a qualitative explanation. A concise answer candidate can
improve token-overlap F1 for slot questions, but remove essential modifiers for
``why'' or ``how'' questions. These failures are not simply retrieval errors.
They are validity-boundary errors: the system has stored what worked without
storing the condition under which it worked.

This paper argues for a different principle:

\begin{quote}
\emph{Experience memory should store not only guidance, but also the
conditions under which the guidance is expected to be beneficial and the
boundaries under which it must be suppressed.}
\end{quote}

We instantiate this principle in \method{} (Conditional Applicability and
Validated Experience for Memory Search), a training-free controller for
memory-search QA. Instead of always injecting retrieved experience, \method{}
casts experience reuse as an intervention-selection problem. Each candidate
experience is represented as
\[
e = (g, a, b, \hat{u}),
\]
where $g$ is the guidance, $a$ is an applicability predicate over the current
question and memory substrate, $b$ is a set of failure-boundary checks, and
$\hat{u}$ is an empirical utility estimate from previous interventions or
diagnostics. At inference time, a candidate is applied only if
\[
a(q, s, m)=1,\quad b(q, s, m)=0,\quad \hat{u}(e)>0.
\]
Otherwise, the controller abstains and preserves the base agent's behavior.

Our implementation targets the same GPT-4o-mini setting used by the local
GAM/R2Mem API harness. The system combines three validated components:

\begin{enumerate}[leftmargin=*]
  \item \textbf{Substrate-aware applicability.} The agent profiles whether the
  memory is episodic dialogue, fact-pack QA, or narrative text. Date/count
  evidence interventions are reused only on substrates where their validation
  predicts positive utility.
  \item \textbf{Answer-contract validation.} The controller parses the answer
  contract implied by the question and only accepts candidate answers whose
  granularity matches the contract, e.g., short slot answers for direct
  ``what/where/who/when'' questions.
  \item \textbf{Failure-boundary suppression.} Boundary checks block known
  harmful transformations: date-to-year truncation, dropped essential
  modifiers, off-slot attributes, lost temporal relations, unsupported count
  changes, and over-compression of explanatory answers.
\end{enumerate}

On LoCoMo, HotpotQA, and NarrativeQA, this conservative intervention policy
outperforms both GAM and R2Mem under the same GPT-4o-mini backbone. The result
is especially informative because the method often abstains: HotpotQA already
satisfies the answer-contract checks after the base CAVE profile, while
NarrativeQA improves from only 16 accepted answer-level interventions out of
300 questions. The gains therefore do not come from injecting more experience;
they come from deciding when experience should not be trusted.

\paragraph{Contributions.}
This draft makes three claims.

\begin{itemize}[leftmargin=*]
  \item We identify \emph{validity-boundary failure} as a core weakness of
  flat experience banks for memory search: retrieved experience can be
  plausible and relevant but conditionally harmful.
  \item We propose \method{}, a training-free framework that augments
  experience reuse with applicability conditions, failure boundaries, and
  utility-bounded intervention selection.
  \item We provide matched full/common GPT-4o-mini evaluations showing that
  \method{} improves over GAM and R2Mem on LoCoMo, HotpotQA, and NarrativeQA.
\end{itemize}

\section{Related Work}

\paragraph{Memory-search agents.}
GAM-style agents decompose long-memory QA into repeated planning, retrieval,
integration, and answer generation. This gives the model a structured search
interface but does not by itself learn from prior search failures. R2Mem adds
reflective experience memory: it scores past trajectories, distills good and
bad experience, and retrieves experience during future search. Our method keeps
the same motivation but changes the decision problem. Instead of asking only
which experience is semantically relevant, we ask whether the experience is
valid under the current memory substrate and answer contract.

\paragraph{Procedural and reflective memory.}
Many recent agent-memory systems store procedural knowledge, reflections, or
skills from prior episodes. These systems differ in how they summarize,
retrieve, update, or organize experience, but they often assume that a retrieved
experience item is usable once selected. \method{} is complementary: it can be
placed after any retrieval policy as a validity layer that decides whether the
candidate should actually intervene.

\paragraph{Credit assignment and negative transfer.}
Experience extracted from a successful trajectory can be spurious: the
trajectory may have succeeded despite a reflected rule rather than because of
it. Conversely, a rule that is useful in one dataset can harm another. This
paper focuses on a lightweight, practical version of this issue: rather than
training a new causal model, we attach observed utility and explicit failure
boundaries to candidate interventions and let the controller abstain when the
validity condition is uncertain.

\section{Problem Setup}

Let $q$ be a question, $M$ a memory store, and $A_0$ a base memory-search agent.
The agent produces a trajectory
\[
\tau = (p_1, r_1, z_1, \ldots, p_T, r_T, z_T, y),
\]
where $p_t$ is a search plan, $r_t$ retrieved evidence, $z_t$ the integrated
working memory, and $y$ the final answer. An experience-augmented agent also
has a bank $E=\{e_i\}$ and a reuse policy $\pi(e\mid q,z_t,M)$.

Flat reuse policies optimize relevance:
\[
e_{1:K} = \operatorname{TopK}_{e\in E}\; \score(e, q, z_t).
\]
The selected experience is then placed in the prompt. This formulation cannot
distinguish three cases:

\begin{enumerate}[leftmargin=*]
  \item \textbf{Applicable and beneficial:} the experience changes the search
  or answer in a way that improves the metric.
  \item \textbf{Relevant but inapplicable:} the experience discusses the same
  phenomenon but assumes a different substrate, answer type, or evidence
  availability.
  \item \textbf{Plausible but spurious:} the experience was extracted from a
  successful trajectory but was not the cause of success.
\end{enumerate}

We therefore evaluate an experience by intervention utility:
\[
\Delta(e;q,M) =
R(A_0(q,M,e)) - R(A_0(q,M,\varnothing)),
\]
where $R$ is the task reward or evaluation score. Since exact counterfactual
evaluation is expensive at inference time, \method{} uses cached diagnostics
and deterministic boundaries to estimate whether $\Delta$ is likely positive.

\section{\method{}}

\method{} is a controller wrapped around a base memory-search agent. It does
not train the backbone model. It changes only whether candidate experience
interventions are allowed to affect the search summary or final answer.

\subsection{Experience Representation}

An experience item is represented as
\[
e = (g, a, b, \hat{u}, c),
\]
where:

\begin{itemize}[leftmargin=*]
  \item $g$ is the natural-language guidance or candidate transformation.
  \item $a(q,s,m)$ is an applicability predicate over the question, current
  search state, and memory profile.
  \item $b(q,s,m,g)$ is a boundary predicate that returns true when a known
  harmful condition is detected.
  \item $\hat{u}$ is the observed utility of the intervention on matched
  development diagnostics.
  \item $c$ stores provenance, such as which candidate system proposed the
  intervention and which failure pattern motivated it.
\end{itemize}

The controller applies $e$ only when
\[
\valid(e,q,s,m) =
\mathbb{1}[a(q,s,m)=1]
\mathbb{1}[b(q,s,m,g)=0]
\mathbb{1}[\hat{u}>0].
\]
If no candidate is valid, the output remains unchanged.

\subsection{Memory-Substrate Profiling}

Experience transfer depends on how memory is represented. \method{} first
classifies the memory substrate into coarse types:

\begin{itemize}[leftmargin=*]
  \item \textbf{episodic dialogue:} temporally indexed conversations, common in
  LoCoMo;
  \item \textbf{fact pack:} document collections with localized factual
  support, common in HotpotQA;
  \item \textbf{narrative:} long story summaries or plots, common in
  NarrativeQA.
\end{itemize}

Date/count anchors are treated as candidate experience interventions. They are
validated for episodic dialogue and fact-pack memory, but suppressed on
narrative memory unless the answer contract provides a narrow slot condition.
This prevents the temporal/count heuristics that help LoCoMo from being
unconditionally transferred to NarrativeQA.

\subsection{Answer-Contract Validation}

The question defines an answer contract. Direct slot questions such as
``where'', ``who'', ``what year'', and ``how many'' usually reward concise
answers. Explanation questions such as ``why'' and ``how did'' often require an
event or reason phrase and can be harmed by over-shortening. \method{} uses the
contract to validate candidate final answers:

\begin{itemize}[leftmargin=*]
  \item For direct slot questions, a shorter candidate can replace the base
  answer only if it is contained in, or is a strict token subset of, the base
  answer or is independently agreed on by multiple candidate paths.
  \item For explanation questions, compression is allowed only when the
  candidate is a multi-token event or reason phrase contained in the base
  answer.
  \item Date answers cannot be reduced from a full date to a bare year unless
  the question explicitly asks for a year.
\end{itemize}

\subsection{Failure-Boundary Suppression}

The controller rejects candidate interventions that cross known negative
boundaries. In the current implementation these include:

\begin{itemize}[leftmargin=*]
  \item replacing a full date with a year for a ``when'' question;
  \item introducing a year not present in either the question or base answer;
  \item dropping essential modifiers such as ``blessed by'';
  \item selecting an off-slot attribute from GAM when the question asks for a
  different property;
  \item losing temporal relations, list components, quoted phrases, reasons,
  or feelings that are required by the answer contract;
  \item accepting broad raw-localizer or count/date overrides whose measured
  utility was negative in diagnostics.
\end{itemize}

The key design choice is conservative abstention. A candidate with uncertain
validity is not applied.

\subsection{Algorithm}

\begin{enumerate}[leftmargin=*]
  \item Run the base memory-search agent and collect its answer, working
  memory, and optional candidate outputs from validated auxiliary paths.
  \item Profile the memory substrate and infer the answer contract from the
  question.
  \item Enumerate candidate interventions from typed evidence modules,
  baseline candidates, and answer-level arbitration.
  \item For each candidate, evaluate applicability, boundary checks, and
  observed utility.
  \item Apply the highest-priority valid intervention; otherwise return the
  base output unchanged.
\end{enumerate}

\section{Implementation}

All experiments use the local API-only GAM/R2Mem harness. The harness reuses
the upstream memory agent, research agent, prompts, dataset loaders, and
metrics. It replaces hardware-dependent components with API-compatible
equivalents: OpenAI \texttt{text-embedding-3-small} for dense retrieval and a
pure-Python Okapi BM25 implementation for lexical retrieval.

\paragraph{LoCoMo.}
The LoCoMo implementation composes several validated answer-level
interventions. ELAA performs answer-level arbitration over cached candidate
summaries. VBR applies validity-boundary recalibration. Slot rescue accepts
only reason, feeling, and quote replacements after diagnostics showed broader
slot replacements regressed. The final utility-bounded contract overlay applies
only candidate replacements with positive measured utility and explicit answer
contracts.

\paragraph{HotpotQA and NarrativeQA.}
For cross-dataset evaluation, the base candidate is the CAVE-profile output.
An offline x-dataset overlay considers GAM, R2Mem, and date/count candidates
when available, but only applies replacements that pass the answer-contract and
failure-boundary checks. The overlay makes no new API calls.

\paragraph{Training and compute.}
No model is trained. All reported numbers use GPT-4o-mini as the backbone.
This isolates the effect of the experience controller from backbone scale.

\section{Experiments}

\subsection{Datasets and Metrics}

We evaluate on three memory-search QA benchmarks:

\begin{itemize}[leftmargin=*]
  \item \textbf{LoCoMo:} nine common held-out conversations, excluding
  \texttt{conv-26}, which is used as the LoCoMo demonstration source.
  \item \textbf{HotpotQA:} the \texttt{eval\_400} split used by the local
  GAM/R2Mem harness, 128 questions.
  \item \textbf{NarrativeQA:} the seed-42 300-question subset used by the
  R2M-style setting.
\end{itemize}

LoCoMo is scored with token F1 and BLEU-1 by question category. HotpotQA and
NarrativeQA use maximum token F1 over gold answers. The main table reports F1.

\subsection{Baselines}

We compare against:

\begin{itemize}[leftmargin=*]
  \item \textbf{GAM:} the base planning-search-integrate-reflect agent.
  \item \textbf{R2Mem:} reflective experience retrieval using faithful
  experience banks and the same GPT-4o-mini backbone in our API harness.
\end{itemize}

For R2Mem, stale top-level summaries were rebuilt from per-sample
\texttt{qa\_results.json} files. Missing HotpotQA and NarrativeQA samples were
rerun before scoring the full/common splits.

\subsection{Main Results}

\begin{table}[t]
\centering
\small
\begin{tabular}{lrrrrrr}
\toprule
Dataset & $N$ & GAM & R2Mem & \method{} & vs. GAM & vs. R2Mem \\
\midrule
LoCoMo common no26 & 1388 & 54.50 & 53.88 & 60.35 & +5.85 & +6.47 \\
HotpotQA eval\_400 & 128 & 64.07 & 62.74 & 66.24 & +2.16 & +3.49 \\
NarrativeQA 300 & 300 & 38.66 & 37.24 & 39.37 & +0.71 & +2.12 \\
\bottomrule
\end{tabular}
\caption{Full/common GPT-4o-mini comparison. Values are token F1.}
\label{tab:main}
\end{table}

\method{} outperforms both baselines on all three datasets. The largest gain
appears on LoCoMo, where temporal state, exact-slot, and answer-contract
failures are frequent. HotpotQA improves even though the final x-dataset
overlay abstains, indicating that the CAVE-profile base path already provides
the benefit. NarrativeQA shows a smaller but consistent gain: only 16 of 300
answers are changed, yet the method exceeds both GAM and R2Mem.

\subsection{LoCoMo Category Results}

\begin{table}[t]
\centering
\small
\begin{tabular}{lrrrr}
\toprule
Method & Multi-hop & Temporal & Open-domain & Single-hop \\
\midrule
GAM & 42.32 & 60.98 & 30.62 & 58.64 \\
R2Mem & 44.07 & 58.08 & 30.94 & 57.98 \\
\method{} & 52.26 & 65.21 & 40.77 & 63.29 \\
\bottomrule
\end{tabular}
\caption{LoCoMo category F1 on the common no26 split.}
\label{tab:locomo_category}
\end{table}

The gains are not isolated to temporal questions. The answer-contract overlay
also improves multi-hop and open-domain questions, suggesting that many
failures occur after evidence is already partially present but the final answer
violates the expected granularity.

\section{Analysis}

\paragraph{Abstention is important.}
On HotpotQA, the cross-dataset overlay makes zero final-answer changes because
the base answers already satisfy the conservative answer-contract checks. This
is a success case for the framework: a controller that must always intervene
would risk negative transfer, while \method{} preserves the stronger base path.

\paragraph{Small numbers of valid interventions can matter.}
On NarrativeQA, the overlay changes only 16 answers. The accepted changes are
primarily shorter contract-satisfying answers supported by GAM or by agreement
between candidate paths. This raises F1 from 38.73 to 39.37.

\paragraph{Utility-bounded overlays reduce negative transfer.}
On LoCoMo, the final utility overlay makes 57 positive changes relative to the
previous VBR+slot-rescue checkpoint and zero negative changes under the cached
candidate diagnostics. Broad candidate auditing, broad raw-localizer usage, and
broad count/date overrides were excluded because diagnostics showed negative
utility.

\section{Limitations}

This draft has several limitations that should be addressed before submission.

\begin{itemize}[leftmargin=*]
  \item The current experiments use GPT-4o-mini only. Qwen 3B/7B/14B code is
  needed to evaluate whether the principle transfers across local backbones.
  \item The API harness reproduces the GAM/R2Mem mechanism. It replaces
  BGE-M3 with API-friendly dense retrieval, and replaces the Lucene-based
  sparse retriever with pure-Python BM25.
  These deviations must be reported.
  \item The current validation is utility- and boundary-based, not a full
  randomized causal validation of every natural-language experience.
  \item Some LoCoMo contract rules are derived from failure analysis on the
  held-out run. Before a final paper, they should be converted into cleaner
  rule families or validated on an additional split.
  \item Statistical uncertainty is not yet reported. Bootstrap confidence
  intervals or paired significance tests should be added.
\end{itemize}

\section{Conclusion}

Experience memory for agents should not be treated as an append-only bank of
generally useful advice. Reflections can be relevant but invalid, useful in one
memory substrate but harmful in another, or beneficial only for a narrow answer
contract. \method{} operationalizes this view by attaching applicability,
failure boundaries, and utility estimates to candidate experience
interventions. In GPT-4o-mini memory-search experiments, this conservative
validated reuse policy outperforms GAM and R2Mem on LoCoMo, HotpotQA, and
NarrativeQA. These results support a more general research direction:
self-improving agents should learn not only what worked, but also when it
should be trusted.

\section*{Reproducibility Pointers}

The main result files are:

\begin{itemize}[leftmargin=*]
  \item \path{r2m-api/results/gam-locomo-4omini-full-no26/summary.json}
  \item \path{r2m-api/results/r2mem-locomo-4omini-full/summary.json}
  \item \path{memory_directions/results/locomo-full-utility-overlay-v11-contract-expanded/summary.json}
  \item \path{r2m-api/results/gam-hotpot400/summary.json}
  \item \path{r2m-api/results/r2mem-hotpot400-full/summary.json}
  \item \path{memory_directions/results/hotpotqa-eval400-v11-xdataset-fullr2m/summary.json}
  \item \path{r2m-api/results/gam-nqa300/summary.json}
  \item \path{r2m-api/results/r2mem-nqa300-full/summary.json}
  \item \path{memory_directions/results/narrativeqa300-v11-xdataset-fullr2m/summary.json}
\end{itemize}

\bibliographystyle{plain}
% TODO: Add verified BibTeX entries before submission.
% \bibliography{references}

\end{document}
